\documentclass[11pt,aps,preprint,nofootinbib,superscriptaddress,tightenlines]{revtex4-2}
\usepackage{amsmath,amssymb,bm,mathtools}
\usepackage{xcolor}
\definecolor{darkblue}{RGB}{0,0,160}
\usepackage[colorlinks=true,linkcolor=darkblue,citecolor=darkblue,urlcolor=darkblue]{hyperref}

\def\st{\begin{equation}}
\def\stp{\end{equation}}
\def\Sect#1{Sect.~\ref{#1}}
\newcommand{\avg}[1]{\langle #1\rangle}
\let\oldparagraph\paragraph
\renewcommand{\paragraph}{\par\bigskip\oldparagraph}

\begin{document}

\title{Plan: speed, measurements, and next steps for stochastic SPH}
\author{D.~Teaney}
\affiliation{Department of Physics and Astronomy, Stony Brook University,
Stony Brook, New York 11794, USA}
\date{\today}

\begin{abstract}
The full Hamiltonian stochastic runs work (\texttt{fullruns.tex}).  This note
plans the next stage: where the time goes and how to cut it, running at
NERSC, measuring correlators in thermalized ensembles, guarding against
superluminal proposals, parallelizing the code, and the physics to do once
all of this works.
\end{abstract}

\maketitle

%=============================================================================
\section{Where the time goes}
\label{sec:time}

\paragraph{Profile.}  One time step of the coupled run (stochastic half step,
leapfrog, stochastic half step), single thread, measured while other runs
shared the machine, so absolute times are pessimistic:
\begin{center}
\begin{tabular}{lrr}
\hline
 & $h/a=2$, $N=672$ & $h/a=2.5$, $N=1020$ \\ \hline
one stochastic half step & 20 ms & 37 ms \\
\quad rebuild stars (neighbours, comoving weights, $4\times4$ solve) & 8 & 15 \\
\quad gradient correction (stars, fits, $g_a$ twice) & 3--15 & 5 \\
\quad stress sweep & 3 & 6 \\
\quad residual sweep (single vote) & 3 & 7 \\
\quad residual sweep (per-bond vote) & 8 & 14 \\
\quad smoothing, thermodynamics, particle energies & 4 & 5 \\
leapfrog (ideal SPH) & 5 & 12 \\ \hline
one time step & $\simeq45$ ms & $\simeq85$ ms \\
per particle and step & $\simeq65\ \mu$s & $\simeq85\ \mu$s \\
\hline
\end{tabular}
\end{center}
About $90\%$ of the time is in the stochastic step, and within it the
Metropolis sweeps are only about a third.  Most of the rest is bookkeeping:
rebuilding the stars and the fits from scratch twice per step, with many
small allocations.  One equation-of-state inversion costs $0.8\ \mu$s.

\paragraph{Cost of a run.}  The cost is
\st
\text{wall time}\simeq c\times N\times N_{\rm steps},\qquad c\simeq65\ \mu{\rm s}\ (\text{now, one thread}).
\stp
The present test box ($N=672$, $t=100$, $11{,}500$ steps) takes about 8 minutes
per realization on one thread.  A box of side $100h$ at $h/a=2$ has $N\simeq46{,}000$;
to $t=100$ it would take about $10$ hours per realization on one thread now.

\paragraph{Options, cheapest first.}
\begin{enumerate}
\item \emph{Merge the two stochastic half steps.}  The second half step of one
  time step and the first of the next act at the same positions.  Doing one
  stochastic step of length $\Delta t$ between leapfrogs (and half steps only at
  the ends of the run or at measurements) halves the stochastic cost.  The
  proposals grow by $\sqrt2$, so the acceptance drops somewhat.  Saving: about
  $45\%$.
\item \emph{Build the stars once per step and reuse the neighbour lists.}
  Keep a neighbour list with a skin and rebuild it only when a particle has moved
  more than half the skin; recompute only the weights.  Share the stars between
  the bond weights and the gradient fit.  Saving: most of the star cost.
\item \emph{Remove allocations in the hot loops:} preallocated buffers,
  static $2\times2$ fits in closed form instead of \texttt{pinv}, arrays of
  structures replaced by flat arrays.  Typical gain in Julia: a factor 2--5.
\item \emph{Single vote} instead of the per-bond vote: half the residual cost
  and higher acceptance.
\item \emph{Threads} (\Sect{sec:parallel}) and \emph{ensembles} (\Sect{sec:nersc}).
\end{enumerate}
Items 1--4 together should bring $c$ to roughly $10\ \mu$s per particle and step on
one thread, a factor $5$--$7$.

%=============================================================================
\section{Running at NERSC}
\label{sec:nersc}

The CPU nodes of Perlmutter have 128 cores (two 64-core AMD EPYC sockets).
Julia is available as a module, and the package environment of the repository
(\texttt{Project.toml}, \texttt{Manifest.toml}) can be instantiated in the user
area; PyPlot is not needed on the compute nodes (plots are made afterwards).
\begin{itemize}
\item \emph{Ensembles first.}  Correlators need many independent realizations.
  These are embarrassingly parallel: one realization per group of threads,
  many groups per node, many nodes through a job array.  This needs no change
  to the code beyond writing each realization's output to its own file.
\item \emph{Threads within a realization} for large boxes (\Sect{sec:parallel}).
\item \emph{Domain decomposition} across nodes only if single realizations no
  longer fit on one node; not needed soon.
\end{itemize}
Practical steps: install and precompile the environment, a small driver
script that reads a parameter file and a seed, writes HDF5 (or JLD2) output,
and can restart from a checkpoint; a Slurm job-array template.

%=============================================================================
\section{Measurements in thermalized ensembles}
\label{sec:measure}

\paragraph{Start in equilibrium.}  The slowest modes relax on the diffusion
time of the box, $1/(\nu k_{\rm min}^2)$ with $\nu=\eta/w$ and $k_{\rm min}=2\pi/L$.  For the
present parameters $\nu=2.2\times10^{-3}$; this is about $4$ time units for $L=0.6$
but about $300$ for $L=5$.  Waiting for thermalization from rest is therefore
not an option for large boxes.  Instead, start from an equilibrium draw:
momenta drawn from the single-particle covariance (\texttt{stochastic.tex}, the
moving fluid), with the total momentum fixed, and a short burn-in for the
particle positions and the lattice modes.

\paragraph{What to record.}  For a fixed set of wave vectors (all $\bm k$ with
$|\bm k|$ below a cut, say $kh<1.5$), record at regular intervals the particle sums
\st
\bm M_{\bm k}(t)=\sum_a\bm m_a(t)\,e^{-i\bm k\cdot\bm x_a(t)},\qquad
S_{\bm k}(t)=\sum_as_a(t)\,e^{-i\bm k\cdot\bm x_a(t)},\qquad
E_{\bm k}(t)=\sum_ae_a(t)\,e^{-i\bm k\cdot\bm x_a(t)} ,
\stp
the Fourier transforms of the momentum, entropy and energy densities.  For
$kh\ll1$ they coincide with the transforms of the smoothed fields.  Split the
momentum into its longitudinal and transverse parts,
\st
M^L_{\bm k}=\hat{\bm k}\cdot\bm M_{\bm k},\qquad M^T_{\bm k}=\hat{\bm k}\times\bm M_{\bm k} .
\stp
This is a cost of order $N\times N_k$ per record, small if recorded every few
steps.  Store the time series per realization (HDF5), with the energies, the
acceptance and the parameters.

\paragraph{The transverse correlator.}  For example, with $\bm k$ along $x$,
\st
C_T(k,\tau)=\frac{1}{V}\,\avg{M_y(t+\tau,\bm k)\,M_y^*(t,\bm k)}
\stp
(the asked-for $\avg{m_y(t,k)\,m_y(0,-k)}$, since $M_{-\bm k}=M_{\bm k}^*$).  Estimate it
by averaging over the start time $t$ in the stationary state (by FFT of the
time series, or directly over lags), over all $\bm k$ of the same length (with
$M_y$ replaced by $M^T$), and over realizations.  Errors from the spread over
realizations, or by jackknife over time blocks longer than the correlation
time.  Prediction from linearized hydrodynamics:
\st
C_T(k,\tau)=\chi_T\,e^{-\nu k^2|\tau|},\qquad \chi_T=wT ,
\stp
so the equal-time value checks the statics and the decay rate measures the
shear viscosity at that $k$, including its renormalization by the
fluctuations.

\paragraph{The longitudinal correlator.}  $C_L(k,\tau)$ from $M^L$: two damped sound
modes,
\st
C_L(k,\tau)\simeq\chi_L\,e^{-\frac12\Gamma_sk^2|\tau|}\cos(c_sk\tau),\qquad
\Gamma_s=\frac{\tfrac43\eta+\zeta}{w}\quad(\text{3D fluid, planar flow}),
\stp
plus a small non-oscillating correction.  Together with $\avg{E_{\bm k}E_{\bm k}^*}$,
$\avg{S_{\bm k}S_{\bm k}^*}$ and their cross correlators this gives the full set of
static susceptibilities and the sound speed and attenuation.

\paragraph{Errors.}  Two levels.  Realizations are independent, so the mean
and the standard error over realizations are the most robust estimate.  Within
a realization the time series is correlated over the decay time $1/(\nu k^2)$,
which is long at small $k$; binning the series into blocks longer than this
(and measuring the integrated autocorrelation time) gives honest errors from
few long runs.  Derived quantities (a fitted $\nu(k)$, ratios of correlators)
get errors by jackknife over realizations or time bins.  Julia packages:
\texttt{Statistics} (standard library), \texttt{BinningAnalysis.jl} (binning,
jackknife, autocorrelation times for Monte Carlo time series),
\texttt{MCMCDiagnosticTools.jl} (effective sample size, agreement between
independent chains), \texttt{LsqFit.jl} (fits with parameter covariances),
\texttt{FFTW.jl} (correlators by FFT), \texttt{Bootstrap.jl}.

\paragraph{Tools to build.}
\begin{itemize}
\item An observer hook in the time loop: a list of measurements, each called
  every $n$ steps, writing to the output file.
\item The equilibrium initial state.
\item A post-processing module: read the files, correlators by FFT, averages
  over $|\bm k|$ shells and realizations, errors, fits of the decay rates, and
  comparison with the predictions above.
\item Spectra $\avg{|M^{L,T}_{\bm k}|^2}$ versus $k$ as the first check of every run.
\end{itemize}

%=============================================================================
\section{Guarding against superluminal proposals}
\label{sec:guard}

A proposal can push a particle's momentum beyond what its energy allows.  The
inversion of the equation of state then fails (or gives $|v|\geq1$).  The
fix is simple and exact: such a state has zero weight under $e^{\mathbb S}$, so
the move is rejected.  In practice
\begin{itemize}
\item in the vote, a failed or superluminal inversion for any member of the
  star counts as $\Delta\mathbb S=-\infty$, and the whole move is rejected;
\item in the deterministic residual ($\lambda=\infty$) and in the back map, a
  superluminal result is clipped or the step is subdivided, and counted;
\item counters for these events are written to the output.
\end{itemize}
Detailed balance is unaffected, since the reverse of a rejected move is never
proposed from an allowed state.  This is a few lines in \texttt{ask\_star!} and
should be done before the long production runs.

%=============================================================================
\section{Parallelization}
\label{sec:parallel}

\paragraph{Do not start from scratch.}  The algorithm is now validated, and
the profile shows that the cost is dominated by avoidable bookkeeping, not by
the structure of the method.  A rewrite in another language would cost months
and the validation would have to be repeated.  The plan is to refactor the
hot paths in Julia, keeping the present code as the reference that every
refactored routine must reproduce (same random streams, same results).

\paragraph{Steps.}
\begin{enumerate}
\item A benchmark script with the per-part timing above, run before and after
  every change; Julia's profiler and allocation tracker on the hot loops.
\item The single-thread items of \Sect{sec:time}.
\item Threads: the smoothing, thermodynamics, velocities and leapfrog loops
  are per particle and thread trivially.  The stochastic sweeps already run the
  cells of one colour in parallel; with many cells (large boxes) this scales,
  with few cells (the present test box has 30) it does not, and ensembles are
  the better use of the cores.
\item Later, if needed: domain decomposition with a halo of width $2h$, which the
  checkerboard of cells already respects.
\end{enumerate}
The random streams are per particle, so the results do not depend on the
number of threads; this is the check that each step preserves.

%=============================================================================
\section{Next steps, once this works}
\label{sec:next}

\begin{enumerate}
\item \emph{Equilibrium validation.}  Static susceptibilities, $C_T$ and $C_L$
  versus $k$ and $\tau$, in boxes large enough to resolve a range of $kh\ll1$.
  The decay rates give $\eta$ and $\Gamma_s$ as functions of $k$: their approach to
  the input values as $k\to0$, and the fluctuation correction.
\item \emph{Renormalization of the viscosity.}  In two dimensions the
  fluctuation correction to $\eta$ grows logarithmically at small $k$ (long-time
  tails).  Measure $\eta(k)$ over a decade in $k$ and compare with the
  one-loop prediction, including the cutoff dependence on $h$.
\item \emph{Dependence on the discretization:} $h/a$, $\Delta t$, $\lambda$, and the
  single versus per-bond vote, for the measured transport coefficients.
\item \emph{Energy drift.}  Decide between a global thermostat and a fix of
  the It\^o term of the change of representation (\texttt{fullruns.tex}).
\item \emph{Moving and inhomogeneous backgrounds} with the Hamiltonian step on:
  a uniformly moving fluid (Lorentz covariance of the correlators), then
  shear flow and an expanding (Bjorken-like) background.
\item \emph{Physics targets:} long-time tails and the renormalization of the
  transport coefficients in relativistic fluids, fluctuations near a critical
  point, and the density-frame formulation that motivated the scheme.
\end{enumerate}

%=============================================================================
\section{Status}
\label{sec:status}

On the branch \texttt{nextstage}:
\begin{itemize}
\item \emph{Superluminal guard}: done (\texttt{test/teststochastic\_guard.jl}).  A
  failed or superluminal inversion rejects the move; an unphysical intermediate
  state of the bond split falls back to the even split; the gradient back map
  falls back to the plain map; the deterministic residual subdivides its step.
  Counters are reported.
\item \emph{Merged stochastic half steps}: done in the new driver
  (\texttt{examples/equilibrium.jl}); about half the cost per step.
  The residual acceptance drops from $0.88$ to $0.77$ (single vote).
\item \emph{Equilibrium initial state, recorder, correlators}: first version
  (\texttt{src/measure.jl}, \texttt{examples/stochastic\_correlators.jl}).
\item \emph{Static solves} in the comoving weights and the gradient fit: done,
  no measurable speed-up; a singular moment matrix (very fast flow, $\gamma\gtrsim2$,
  when too few bonds remain inside the support along the flow) now falls back
  instead of failing.  The comoving stencil at large $\gamma$ needs a closer look.
\item \emph{Next}: a proper profile (sampling profiler and allocation tracking)
  of \texttt{build\_stars} and the gradient correction.
\end{itemize}

%=============================================================================
\section{Order of work}
\label{sec:order}

\begin{enumerate}
\item The superluminal guard (\Sect{sec:guard}).  Small; do first.
\item Merge the stochastic half steps; build the stars once per step.  Check
  that the results are unchanged in distribution.
\item Equilibrium initial state, observer hook, HDF5 output, correlator
  post-processing; first measurement of $C_T$ in the present box.
\item Allocation and fit clean-up; threads; benchmark.
\item NERSC: environment, driver, job-array template; first ensemble in a
  larger box.
\item The physics of \Sect{sec:next}.
\end{enumerate}

\end{document}
